\section{GBC for the initial condition problem}\label{m:sec:gbc}

Between the forward solves of Steps 1 and 4 of Algorithm~\ref{m:alg:gbc} sits a learning problem that is finite dimensional whatever $N$ is, because what the generator conditions on is a summary of the data and not the data.

Fix the model \eqref{m:eq:model} to \eqref{m:eq:prior} and a summary map $S_N:(\R^W\times\mathcal X)^N\to\R^{p_N}$. Let $U$ be a base random variable independent of everything.

\begin{algorithm}[GBC: solve for the initial condition, simulate forward]\label{m:alg:gbc}\leavevmode
\begin{enumerate}[nosep]
\item Simulate $M$ training pairs: $\theta^{(m)}\sim\Pi$, then $Z^{(m)}$ from \eqref{m:eq:model} with $\Gcal(\theta^{(m)})$ computed by a forward solve. Compute $S^{(m)}=S_N(Z^{(m)})$.
\item Train a conditional generator $\widehat G:\R^{p_N}\times\R^q\to L^2_0(\Omega)$ on $\{(S^{(m)},\theta^{(m)})\}$ so that $\Law(\widehat G(s,U))$ approximates $\Law(\theta\mid S_N(Z)=s)$. In coefficient form, $\widehat G(s,U)=\sum_j\widehat Q_j(U_j\mid s)e_j$ with $\widehat Q_j(\cdot\mid s)$ a learned conditional quantile function.
\item Define the GBC posterior $\Pihat:=\Law\big(\widehat G(S_N(Z^{(N)}),U)\big)$ conditional on the observed $Z^{(N)}$, and its mean $\hat\theta_N^{\rm G}$.
\item Simulate forward: for $\theta\sim\Pihat$, solve the PDE to obtain $u_\theta$; report $\Law(u_\theta(t)\mid Z^{(N)})$, functionals $\Phi(u_\theta)$, and credible bands.
\end{enumerate}
\end{algorithm}

Step 2 writes the generator in coefficient form, with one learned quantile function for each mode. The generator of Definition~\ref{m:def:gen} instead mixes the coordinates of the base noise through a full matrix, so it has this form only when that matrix is diagonal; the first of the four points in that definition explains why nothing downstream depends on the difference.

Three sample sizes appear and do different work. $N$ is the number of observations, and is the only one the asymptotics of Sections~\ref{m:sec:jw} to~\ref{m:sec:main} take to infinity. $M$ is the simulation budget, chosen by the user, and it is the quantity the generator theorems bound (Theorems~\ref{m:thm:qnnA} and~\ref{m:thm:raw}). The number of draws from $\Pihat$ in Steps 3 and 4 carries no statistical content, since $\Pihat$ is available as a map and can be sampled to any accuracy for the price of forward passes.

Steps 1 and 2 never see $Z^{(N)}$. They can be run before the data exist, and their cost, $M$ forward solves and one training run, is paid once; a new data set then costs one evaluation of $S_N$, one forward pass of $\widehat G$ per draw, and whatever forward solves are asked of Step 4. That is what amortization means here, and it is a claim about cost and not about accuracy.

Step 1 already constrains the summary map. It is applied to simulated data, so it must be a statistic: a function of $Z$ alone, with no dependence on the $\theta$ that produced it. This is why the one step summary of Section~\ref{m:sec:suff} is centered at a preliminary estimate computed from the data rather than at the truth, and why that estimate is cross fitted.

Write $\PiN=\Pi(\cdot\mid Z^{(N)})$ for the exact posterior and $\Pi(\cdot\mid S_N)$ for the posterior given the summary alone. Two discrepancies arise: $\Pi(\cdot\mid S_N)$ versus $\PiN$, which is sufficiency and the subject of Sections~\ref{m:sec:suff} and~\ref{m:sec:jw}, and $\Pihat$ versus $\Pi(\cdot\mid S_N)$, which is generator fit and enters the transfer theorem of Section~\ref{m:sec:main} through one accuracy assumption.

\section{What is taken, what is added, and what is conditional}\label{m:sec:contrib}

The results of Nickl~\cite{nicklbvm2024} and Konen and Nickl~\cite{konennickl2026} are used as black boxes and are not reproved. We take from them the model \eqref{m:eq:model} with the prior \eqref{m:eq:prior}, Conditions 2A to 2F on the forward map, Theorems~\ref{m:thm:nickl2} and~\ref{m:thm:nickl1}, the centering \eqref{m:eq:thetahat}, the Fredholm decomposition \eqref{m:eq:fredholm}, the exact coverage corollary and the local asymptotic minimax theorem, and the inverse logarithmic barrier of Nickl and Titi~\cite{nickltiti2024}. What follows is new relative to that framework, grouped by kind and in the order of the chapters that prove it.

\subsubsection{The sufficient statistic (Section~\ref{m:sec:suff})} For a linear evolution the adjoint equation driven by the data is exactly sufficient for the initial state, and the GBC quantile map is exactly affine with an explicit scale (Lemma~\ref{m:lem:adjoint}, Proposition~\ref{m:prop:linsuff}). For a nonlinear evolution the one step efficient score statistic in $J_N\asymp N\delta_N^2$ modes is $\sqrt N$ equivalent to Nickl's centering, hence a computable center for the Gaussian limit of the posterior (Definition~\ref{m:def:summary}, Proposition~\ref{m:prop:suff}), and it is asymptotically sufficient for the whole law under Assumption~\ref{m:ass:U} (Theorem~\ref{m:thm:jwbvm}); in the linear model its regularized form \eqref{m:eq:regsummary} is exactly the posterior mean, whatever the linearization point (Lemma~\ref{m:lem:exact}), and for the heat and Stokes equations that form remains $\sqrt N$ equivalent to Nickl's centering when $J_N$ modes are retained under the rescaled prior (Proposition~\ref{m:prop:regheat}). The same statistic then does three further things. It is recursively updatable in the linear case, so particle learning for the initial condition is exact without weight degeneracy (Proposition~\ref{m:prop:storvik}). It identifies the Bernstein von Mises covariance as a change in expected utility in the sense of Polson and Roberts~\cite{polsonroberts1993}, exactly in the linear Gaussian models of the heat and Stokes equations and, for a nonlinear forward map in the setting of Theorem~\ref{m:thm:nickl2} and a fixed number of modes, through the realized information gain of the posterior at the truth, whereas the expected gain is bounded by the information near the center of the prior (Corollary~\ref{m:cor:prreal}, Propositions~\ref{m:prop:prvt} and~\ref{m:prop:prrd}, and Remark~\ref{m:rem:prscope}), and separates the noise level from the initial condition by a quadratic variation argument set in the diffusion framework of Polson and Roberts~\cite{polsonroberts1994} (Propositions~\ref{m:prop:pr} and~\ref{m:prop:qv}). And it is what the summary selection literature selects: under Assumption~\ref{m:ass:U} it attains, asymptotically, the optimum of the four approximate Bayesian computation (ABC) criteria, of the $L^2$ Bayes risk for the exact posteriors under the moment bound~\eqref{m:eq:momentbound} and of the other three in the $J_N$ mode Gaussian projection, which is the limit for heat and Stokes; in the linear case the statistic is the partial least squares (PLS) projection; and in general it has the structure of deep PLS, introduced by Polson, Sokolov and Xu~\cite{polsonsokolovxu2021} and applied to instrumental variable regression by Nareklishvili, Polson and Sokolov~\cite{nareklishvili2023} (Propositions~\ref{m:prop:abcopt}, \ref{m:prop:pls} and~\ref{m:prop:dpls}).

\subsubsection{Sufficiency theorems and bounds (Sections~\ref{m:sec:jw} to~\ref{m:sec:suffbounds})} The posterior mean summary is first moment exact at every $N$ and, under the uniform Bernstein von Mises statement of Assumption~\ref{m:ass:U}, sufficient for the whole law at scale $\sqrt N$ (Theorems~\ref{m:thm:jw} and~\ref{m:thm:jwbvm}). When a summary is not sufficient the gap it leaves is bracketed below by its Rao Blackwell trace, with the operator norm correction of Theorem~\ref{m:thm:gap}, and above by $\sqrt{2\norm{V_S}_{op}\Delta I}$, both computable from the model (Theorem~\ref{m:thm:gap} in the linear Gaussian model, Proposition~\ref{m:prop:gapnl} beyond it, for a summary whose posterior, like the posterior given all of the data, is close at the $\sqrt N$ scale to that of one linear Gaussian model under the prior predictive law). Under the rescaled prior and clauses (U1) to (U3) of Assumption~\ref{m:ass:U} the posterior covariance of that model must be asymptotically the inverse information at the prior mean (Proposition~\ref{m:prop:gapreq}), and in a two coordinate nonlinear model the maximum likelihood coordinate is such a summary, while the first coordinate of the adjoint statistic at the prior mean is not, for any Gaussian model, once $\sqrt Ns_N^2\to\infty$ (Proposition~\ref{m:prop:gaptoy}). And the dimension the generator has to learn in is small under dissipation: the effective dimension of the sufficient statistic is of order $(N\rho^2/\nu T)^{d/(2\alpha+2)}$ as $N\to\infty$ at fixed $\nu T$, decreasing in viscosity (Theorem~\ref{m:thm:jeff}), against a simulation budget that is super exponential on the raw data for a generator required to be uniformly accurate and, conditionally on the hypotheses of Theorem~\ref{m:thm:qnnA}, linear in that effective dimension for the hybrid generator of Theorem~\ref{m:thm:jeff}(ii) trained on the statistic (Theorem~\ref{m:thm:raw}).

\subsubsection{Transfer, generators, and network priors (Sections~\ref{m:sec:main} to~\ref{m:sec:qnn})} Any $\sqrt N$ accurate, concentrating approximation to the posterior inherits the Gaussian limit, the supremum norm trajectory theorem, exact coverage and minimax efficiency (Theorem~\ref{m:thm:main}, whose proof isolates the properties of the posterior that Nickl's trajectory argument uses), and solving forward from the GBC draws agrees asymptotically with simulating the tangent equation from that limit (Corollary~\ref{m:cor:commute}). Which generators qualify is then settled in both directions: conditionally on Assumptions~\ref{m:ass:U} and~\ref{m:ass:u2p} and on a bound for the error of the summary, the identity skip Gaussian head quantile generator meets the accuracy hypothesis when $J_N/M_N\to0$ (Theorem~\ref{m:thm:qnnA}), and no generator required to be uniformly accurate over H\"older quantile maps can meet it with polynomially many simulations once $J_N\to\infty$ (Theorem~\ref{m:thm:qnnB}). Under the Polson Ro\v{c}kov\'a network prior, adaptive contraction transfers to the initial condition through the stability estimate (Theorem~\ref{m:thm:PR}) but the Bernstein von Mises step does not, because its Cameron Martin step has no network analog, and Section~\ref{m:sec:qnn} says what would be needed.

\subsubsection{Localized training (Section~\ref{m:sec:loc})} Training on a proposal of radius $c$ leaves a residual with mean $P\bar\theta$ and, in a basis that diagonalizes the prior and the Gram matrix together, variance $\Sigma_{N,j}(s_j+c^2p_j^2)$, so the mean has to be removed and the variance bias is of first order in the prior fraction (Theorem~\ref{m:thm:locbias}). The three error terms $\beta(c)^2+a_v/c^v+bc^4$ then fix, for the affine head, an optimal radius $c^\star\asymp(DN/Mc_3^2)^{1/6}$, which Figure~\ref{m:fig:radius} checks against measurement (Proposition~\ref{m:prop:radius}). What comes out is a design rule, Algorithm~\ref{m:alg:loc}: localize at $c\in[2,3]$, a calibration for the identity head and not a consequence of \eqref{m:eq:cstar} (Remark~\ref{m:rem:radiushead}), center and rescale the residuals, use the identity head, and do not importance weight.

\subsubsection{Examples and dispersive equations (Sections~\ref{m:sec:heat} to~\ref{m:sec:nondiss})} For the heat equation a matrix Bernstein bound controls the finite $N$ error of the Bernstein von Mises covariance, checked numerically against the exact posterior, with relative discrepancy below one percent in the calibrated runs (Proposition~\ref{m:prop:gram}, Figure~\ref{m:fig:heatgbc}, Table~\ref{m:tab:diag}), and Table~\ref{m:tab:iota} puts the Nickl Titi barrier in numbers. For Navier Stokes this part supplies the Stokes limit, the Taylor Green background, the explicit Fourier coupling \eqref{m:eq:tgcoupling} and the adjoint pipeline, together with the identity \eqref{m:eq:modeinfo} that reconciles Konen Nickl and Nickl Titi and that would give the curvature constant of Nickl and Wang if their remaining conditions were verified (Section~\ref{sec:gradstab}). Nothing in the framework treats dispersive equations, and two results say what changes there: for linear Schr\"odinger the information operator is exactly $T\cdot\Id$, also with a real potential (Section~\ref{m:sec:nondiss}), and for the free equation the transfer theorem holds in a reduced form, without a supremum norm statement (Proposition~\ref{m:prop:sch} and Corollary~\ref{m:cor:schtransfer}), while for cubic NLS about a plane wave the Bogoliubov information operator is a compact perturbation of $T\cdot\Id$ in the defocusing case and, in the focusing band, exponentially amplified along one direction of each unstable block (Proposition~\ref{m:prop:bog}). Condition 2E holds for its solution map by an energy estimate (Proposition~\ref{m:prop:nlsstab}), while Conditions 2C and 2D are not verified.

\subsubsection{Connections (Section~\ref{m:sec:connections})} Least squares Monte Carlo is the same construction, and its two limit convergence theory transfers (Proposition~\ref{m:prop:lsm}, Theorem~\ref{m:thm:lsmgbc}). Three shorter readings follow: the model error endpoint through Carlin, Polson and Stoffer~\cite{cps1992} and Johannes, Polson and Stroud~\cite{jps2009}, Higdon's calibration as forward emulation set against GBC as inverse emulation, and the sixth problem reading of the dissipative dispersive contrast.

\subsubsection{What is unconditional and what is not} Three kinds of hypothesis appear above and they carry different risk. Conditions 1 and 2 of Nickl~\cite{nicklbvm2024} and the constants of Theorem~\ref{m:thm:nickl2} are \emph{verified} for the dissipative equations treated here and are assumed throughout without further comment. Assumption~\ref{m:ass:G} is a hypothesis about the \emph{output of an algorithm}, namely that the trained generator is accurate at the $\sqrt N$ scale and concentrates; it is the interface of the transfer theorem, it is what a practitioner controls by training longer or simulating more, and Section~\ref{m:sec:experiments} measures it. Assumption~\ref{m:ass:U}, and Assumption~\ref{m:ass:u2p} outside the setting of Theorem~\ref{m:thm:u2ploc}, are of the third kind: unproved \emph{mathematical} statements about the framework itself, which we isolate rather than establish. For a nonlinear forward map the second display of Assumption~\ref{m:ass:u2p} fails in general when the training parameters are drawn from the prior (Remark~\ref{m:rem:u2p}). What each result of Chapters~\ref{ch:adjoint}, \ref{ch:generators} and~\ref{ch:examples} needs is stated with it and in the remark that follows it, and the Notes of those chapters name results that are unconditional in the third sense, that is, rest on Conditions 1 and 2, on Assumption~\ref{m:ass:G} where the generator enters, and on nothing that is open. Two results are conditional at the source. Theorem~\ref{m:thm:jwbvm} requires Assumption~\ref{m:ass:U} (and, for the one step summary, the hypotheses of Proposition~\ref{m:prop:suff}), and Theorem~\ref{m:thm:qnnA} requires it as well, together with Assumption~\ref{m:ass:u2p} and a bound on the error of the summary in norms stronger than those in which Proposition~\ref{m:prop:suff} controls it. What is proved from either theorem, or from Assumption~\ref{m:ass:U} directly, is conditional in turn, among them Corollary~\ref{m:cor:jwgbc}, Proposition~\ref{m:prop:abcopt} (which adds the moment bound~\eqref{m:eq:momentbound}), the second assertion of Corollary~\ref{m:cor:commute}, case (c) of Theorem~\ref{m:thm:lsmgbc}, and the statements on the hybrid generator in Theorems~\ref{m:thm:jeff} and~\ref{m:thm:raw}. Assumptions~\ref{m:ass:U} and~\ref{m:ass:u2p} appear as Open Question~\ref{m:op:unif}, together with the moment bound~\eqref{m:eq:momentbound}. A reader who declines to grant them keeps the transfer theorem and everything downstream of it that does not use them, and loses the claim that this particular generator, with linearly many simulations, is one of the approximations the transfer theorem accepts.\looseness=-1

\section{Relation to the computational literature on Bayesian inverse problems}\label{m:sec:related}

\emph{Function space MCMC.} Stuart~\cite{stuart2010} sets the Bayesian inverse problem on function space, and the preconditioned Crank Nicolson sampler of Cotter and coauthors~\cite{cotter2013} exploits that formulation directly (Section~\ref{sec:dimrobust}). This is the natural competitor, and it is the benchmark used in Section~\ref{m:sec:experiments}. Its guarantees are stronger than anything GBC offers, since it targets the exact posterior itself, whereas GBC targets the law of a Gaussian generator fitted to it: away from the Gaussian case that substitution leaves an approximation error which no simulation budget removes. Nickl and Wang~\cite{nicklwang2024} give polynomial time guarantees for Langevin sampling in a related steady state elliptic model (Section~\ref{sec:gradstab}), and Castre and Nickl~\cite{castrenickl2026} have verified gradient stability for the interaction potential of the McKean Vlasov equation, where the unknown is a parameter of the equation and not an initial condition. What function space MCMC does not offer is amortization: a new data set requires a new chain.

\emph{Ensemble and optimization methods.} Ensemble Kalman inversion~\cite{iglesias2013} propagates a finite ensemble through the forward map and applies a Kalman update in ensemble space. It is derivative free and cheap, and it is what operational weather systems run~\cite{houtekamerzhang2016}, but for a nonlinear forward map Ernst, Sprungk and Starkloff~\cite{ernst2015} show that the large ensemble limit is a linear Bayes estimator rather than the posterior, so the ensemble spread is not a posterior credible set. Variational data assimilation, 4DVar, is the corresponding point estimator: Section~\ref{m:sec:suff} notes that the one step summary is computed by exactly the 4DVar machinery, one forward and one adjoint solve, used here to build a statistic rather than to minimize a functional. On the conventional solve counts given in Section~\ref{m:sec:expcost}, one GBC analysis at $4000$ simulations is $400$ four dimensional variational equivalents and $80$ times a $100$ member ensemble filter; that is the price of amortization.

\emph{Laplace with a low rank Hessian.} Bui Thanh and coauthors~\cite{buithanh2013} form the posterior covariance as a low rank update of the prior, using the fact that the data inform only a few directions. Section~\ref{m:sec:resolution} is the statistical counterpart of that observation: the number of informed directions is $J_N$, a small power of $N$. Since Definition~\ref{m:def:gen} is a plug in Gaussian, GBC and this method target the same object, and the honest claim for GBC is amortization and simulation only access to the forward map, not accuracy beyond Laplace. This is made precise in the point of Definition~\ref{m:def:gen} on what is trained and what is not, and confirmed numerically in Section~\ref{m:sec:experiments}.

\emph{Transport maps and measure transport.} El Moselhy and Marzouk~\cite{elmoselhymarzouk2012} construct a deterministic map pushing the prior to the posterior, and Spantini, Bigoni and Marzouk~\cite{spantini2018} show how low dimensional couplings exploit the same effective dimension structure. A quantile generator is a transport map, triangular when its coordinates are generated one after the other, and the generator of Definition~\ref{m:def:gen} is the affine member of that family up to the choice of root of the covariance $\hat C_M$ that it uses: replacing the symmetric root of its definition by the Cholesky factor makes the map triangular and leaves its law unchanged. Theorem~\ref{m:thm:qnnB} says what the nonlinear members cost.

\emph{Simulation based inference.} Neural posterior estimation (NPE)~\cite{greenberg2019,papamakarios2021} and the wider simulation based inference program~\cite{cranmer2020} learn a conditional density or flow from simulated pairs, which is what Algorithm~\ref{m:alg:gbc} does. Three things distinguish the present setting. The summary is not learned but supplied by the adjoint equation, and Lemma~\ref{m:lem:adjoint} and, under Assumption~\ref{m:ass:U}, Theorem~\ref{m:thm:jwbvm} say it is sufficient rather than merely convenient, which is the usual weak point of that literature. The target is a nonparametric object with a known Gaussian limit, so the accuracy required of the generator is a $\sqrt N$ statement rather than a density estimation statement, which is why Theorem~\ref{m:thm:main} can transfer guarantees at all. And the localized design of Algorithm~\ref{m:alg:loc} is the analogue of sequential NPE's proposal refinement, with Theorem~\ref{m:thm:locbias} supplying the correction that sequential NPE applies by importance weighting, a route that degenerates here, as the Navier Stokes study of Polson and Sokolov~\cite{polsonsokolov2026} finds (Section~\ref{m:sec:expns}). Ghattas and Willcox~\cite{ghattaswillcox2021} survey the surrounding territory.

\emph{Operator learning and amortization.} Neural operators~\cite{li2021,kovachki2023} amortize in the other direction: they learn the forward map $\theta\mapsto u_\theta$ once, after which a new coefficient field costs a forward pass, and any inference then runs on top of the surrogate. GBC amortizes the inverse map $Z^{(N)}\mapsto \Pi(\cdot\mid Z^{(N)})$ instead, so a new \emph{data set} costs a forward pass. The two compose, since a neural operator can supply the simulator that generates the training pairs of Algorithm~\ref{m:alg:gbc}, but they are not substitutes, and the guarantees differ in kind: an operator surrogate carries an approximation error in the forward map that has to be propagated through the posterior, whereas Theorem~\ref{m:thm:main} certifies the posterior directly.

The experiments of Section~\ref{m:sec:experiments} cannot quantify the price of the Gaussian approximation: the benchmark against function space MCMC there is linear and Gaussian, which is exactly the case in which the approximation error of a Gaussian head vanishes identically, so the experiment measures the cost of estimating the generator and not the cost of the substitution. Bounding the latter in a genuinely nonlinear problem is left open, and is stated as such in Section~\ref{m:sec:open}.
